\documentclass[conference]{IEEEtran}
\IEEEoverridecommandlockouts

\usepackage{cite}
\usepackage{amsmath,amssymb,amsfonts}
\usepackage{graphicx}
\usepackage{booktabs}
\usepackage{array}
\usepackage{xcolor}
\usepackage{url}
\usepackage{tikz}
\usetikzlibrary{arrows.meta,positioning,shapes.geometric,calc,fit,backgrounds}
\usepackage[colorlinks=true, linkcolor=blue, urlcolor=blue, citecolor=blue]{hyperref}

\tikzset{
  box/.style={rectangle, rounded corners=2pt, draw=black!70, fill=blue!6,
    align=center, minimum height=8mm, inner sep=4pt, font=\footnotesize},
  db/.style={cylinder, shape border rotate=90, aspect=0.25, draw=black!70,
    fill=orange!12, align=center, minimum height=10mm, font=\footnotesize},
  llm/.style={rectangle, rounded corners=2pt, draw=black!70, fill=green!8,
    align=center, minimum height=8mm, inner sep=4pt, font=\footnotesize},
  flow/.style={-{Stealth[length=2mm]}, draw=black!70},
  step/.style={rectangle, rounded corners=2pt, draw=black!70, fill=gray!8,
    align=center, minimum height=7mm, text width=20mm, inner sep=3pt, font=\scriptsize}
}

\begin{document}

\title{HospitAI: A Hospital Analytics Platform Combining Patient Clustering, Clinical Sequence Mining and AI-Generated Interpretation}

\author{\IEEEauthorblockN{Ioan-Iulian Todor}
\IEEEauthorblockA{\textit{Faculty of Automation and Computer Science} \\
\textit{Technical University of Cluj-Napoca}\\
Cluj-Napoca, Romania \\
todor.io.ioan@student.utcluj.ro}
}

\maketitle
\thispagestyle{plain}
\pagestyle{plain}

\begin{abstract}
Many electronic health records (EHR) are generated in the healthcare environment. Unfortunately, a vast majority of this data remains underutilized, as clustering patient data to analyze longitudinal care histories is highly complex.

In this paper, we present HospitAI, a platform structured around a three-stage pipeline: ETL for data ingestion, the K-Prototypes algorithm for mixed-type patient clustering, and a PrefixSpan sequence-mining framework tailored for the medical context. An AI agent also lets users query the database in natural language, without any technical knowledge. The AI is also used for interpreting our results (clusters and sequences), making them easily understandable for both medical professionals and patients.

Analyzing each patient cluster separately and identifying anomalies allows for the detection of unique care pathways, such as oncology, dialysis, and heart failure treatments. Testing the system on a synthetic dataset of 22,284 patients yielded four stable clusters with distinct sequence patterns, successfully representing actual patient categories.

HospitAI converts complex numerical data into accessible text, using LLMs to provide practical, easy-to-read insights into patient categories and care patterns.
\end{abstract}

\vspace{6pt}

\section{Introduction}

Nowadays hospitals contain large amounts of data that are stored in electronic health records (EHRs). With each patient encounter, data about diagnoses, prescribed medications, care plans and other procedures are recorded. On the theoretical level, it would be possible to use such data to find out meaningful patterns about patients and care. In reality, most information remains unused since the tables are too large, untidy, and complicated for an untrained eye.

There are three types of difficulties associated with this kind of data analysis. The first one is caused by the diversity of patients. One patient may be young with a single condition, while another may be older with multiple comorbidities. As a result, calculating a simple mean for the whole set of data hides the differences between the individuals included in it. Another challenge related to patient care is its temporal character. Clinical care is always represented by a series of actions performed during several weeks or months, and their order matters. The last difficulty is interpretation of the results. Even when a model produces a good result, the output is usually a list of numbers or codes that a clinician or a patient cannot interpret quickly.

In this paper, we introduce HospitAI, a web application designed to solve all three of these problems at once. We combined the data transformation, clustering process, sequence mining and the language processing in such a way that the same data of a patient can be viewed through different lenses with the help of one tool.

To solve the first problem, we apply K-Prototypes \cite{huang1998}, an algorithm designed for mixed datasets. The data has both numeric variables (e.g., age of the patient and the number of drugs) and categorical variables (e.g., gender and major care plan), each having a different weight during the clustering process. To solve the second problem, we create a pipeline on top of PrefixSpan algorithm \cite{pei2001} and modify it to make sure it provides useful patient care sequences rather than repetitive ones. To solve the third problem, we integrate LLMs that summarize statistics in human-understandable terms and also help clinicians to query the database faster.

The contributions of this paper are as follows:

\begin{itemize}
  \item Creation of an end-to-end solution that uses mixed data clustering, clinical sequence mining, and generation of AI-assisted explanations, with a clear differentiation between backend analysis and frontend implementation of the web-application.
  \item Creation of a feature-engineering routine that transforms raw data from multiple EHR tables into numeric and categorical features while taking into account the problem of skewed values and missing vitals.
  \item Adaptation of the PrefixSpan algorithm for domain-specific sequence mining with ranking of patterns according to geometric mean lift score and diversity rules. This sharply reduces trivial output and brings clinically relevant pathways to the front.
  \item Use of an LLM as an assisted ``verbalizer'' based on statistical information and a Natural-Language-to-SQL agent to translate natural language queries to SQL.
\end{itemize}

The rest of the paper is organised as follows. Section~\ref{sec:related} reviews related work on care-pathway mining and patient clustering. Section~\ref{sec:system} provides the system overview. Section~\ref{sec:methods} explains the three analytics components. Section~\ref{sec:results} reports the main results, and Section~\ref{sec:conclusion} discusses limitations and future work.
\section{Related Work}
\label{sec:related}

\textbf{Surveys of patient pathway mining.} There exist two recent surveys that provide a good overview of this area. Rosa et al. \cite{rosa2022modelling} provide a survey of the models and techniques used for patient pathways mining. They distinguish between the sequence mining and process mining approaches. Flothow et al. \cite{flothow2023} provide a survey of various techniques applied in the search of patient sequences of healthcare use, distinguishing clustering approaches, which uncover entire patient trajectories, from pattern mining approaches, which are applied to subsequences.

\textbf{Combining clustering with sequence mining.} The most relevant research hybridizes the two concepts. Najjar et al. \cite{najjar2018} suggest a two-step process, whereby they cluster single healthcare visits using K-Prototypes and derive treatment pathways afterwards. It is an important predecessor of our research because it also uses the mixed clustering technique that we employ. Baudrier et al. \cite{baudrier2022} hybridize their care trajectory clustering on patients with kidney cancer with a sequential pattern discovery algorithm and conclude that their clusters differ from each other by actual results like mortality rate. Vogt et al. \cite{vogt2018} discover care pathway clusters for heart-failure patients from sequences in insurance claims with ambulatory healthcare pathways. Our approach is conceptually similar, yet it differs from theirs in the order of steps involved.

\textbf{Sequence mining from EMR with interpretation.} Perer et al. \cite{perer2015} conduct sequence mining from electronic medical records based on frequent event sequences to patient outcomes, with the major emphasis on visual analysis. Their results show that frequent sequences can detect clinically meaningful care patterns unknown to the physician, which motivates our own attention to pattern quality, since frequency alone is not enough.

\textbf{Clustering of patients on EHR data.} Some studies concentrate on analyzing methods of proper patients' clustering. Yang et al. \cite{yang2024} develop a smart-healthcare solution aimed at clustering patients based on different kinds of information contained in their medical records, to reveal potential health risks and disease patterns across patient groups. Aljohani \cite{aljohani2024} evaluates eight different clustering methods for patients' stratification based on EHR and analyzes them using a multi-criteria approach.

\textbf{Process mining as an alternative.} An alternative set of methods analyzes care processes from event logs. Specifically, Murazzano et al. \cite{murazzano2026} use process mining to analyze chronic respiratory patients and observe a large variability in their care processes, pointing to another direction of analysis that addresses the same objective in a different way.

Clustering and sequence mining are both mature domains in healthcare analytics. There have been some research works which try to link them together. However, there remains the need for a unified framework that integrates cluster-scope mining with a rank of clinically meaningful patterns followed by an AI-based explanation component. That's where HospitAI comes in.
\section{System Overview}
\label{sec:system}

\subsection{Architecture}

The HospitAI architecture consists of two components, one that runs online via a web application, and another one which is offline analytics. In the online component, we use a frontend built with Next.js and a backend with FastAPI. The backend implements the layered approach without any bypasses between layers. The layers consist of routers which use API services, service layer which communicates with a repository layer, and repository layer which directly connects to the database. The entire persistence storage consists of a single PostgreSQL database. Figure~\ref{fig:arch} illustrates the overall architecture.

The logic for analytics is separated from the online requests and represents a three-step offline pipeline implemented in a set of Python scripts: ETL cleaning of the initial export data, a clustering step and sequence mining. Each step reads its input from the database, performs its computations and stores the results back into special tables, which the web app then only reads to provide a snappy experience. Above all this sits a single large language model running on top of the OpenAI API, enabling each of the text-generation tasks, including cluster/sequence interpretation and question-answering.

\begin{figure}[t]
  \centering
  \includegraphics[width=0.95\linewidth]{figs/high_level_architecture.png}
  \caption{High-level architecture of the application}
  \label{fig:arch}
  \end{figure}

\subsection{Data}

For training and experimenting, we will use a simulated database created using the open source patient simulator Synthea. Simulated data fits well here because it avoids privacy and legal concerns while keeping a realistic database schema. Our database contains approximately 22,284 patients and standard clinical data tables such as encounters, conditions, medications, procedures, observations and care plans. The ETL process described in Section~\ref{sec:methods} filters our database down to a single hospital and fixed set of providers as well as fixing the dates and calculating some auxiliary information such as patient age and treatment duration.

Both the clinical database and the analysis results are stored in the same database. The clustering algorithm creates a single row for each patient with its clustering result and a list of selected features cached in the row, while sequence mining stores patterns in a separate table along with a mapping of patterns to the patients matching them.
\section{Methods}
\label{sec:methods}

The analytics component of HospitAI consists of an offline pipeline with three stages plus the AI component. Stage one consumes the exports and generates a clean relational database (ETL process). Stage two computes the patient features and clusters patients via the K-Prototypes algorithm. Stage three computes the clinical sequences in each patient group. Above all this, the AI component will take in the preprocessed information and generate readable text output, either as automatically generated cluster/sequence descriptions or as a question-answer session for humans. We explain the details below.

\subsection{Stage 1: Data Ingestion and ETL}

The data flow in HospitAI begins with the comma-separated exports generated by Synthea and ends with a clean PostgreSQL database.

Three key decisions related to reduction and data cleansing influence the results of our analysis. First, we fix one hospital. In particular, the organization which was selected has the highest utilization. Only busy providers are kept. Three hundred providers are ranked by decreasing number of encounters. Then all encounters of the excluded providers are redistributed between other providers with a fixed random seed so that the number of encounters and the provider load remain unchanged. Second, we fix the dates. Some events created by Synthea may have the timestamp set to some date in the future. Therefore, we restrict the timestamp to the current date. We obtain the patient's age by subtracting the birth date from the death date (or from the current date for patients who are still alive). The length of an encounter is estimated using the beginning and end of the encounter time and limited by some reasonable value depending on the encounter's class (for instance, up to two hours for the wellness visit and thirty days for the inpatient encounter). Third, missing values are not mean-imputed, since that would distort the numerical centroids K-Prototypes treats as real measurements. Sparse categorical and vital-sign features instead fall into an explicit ``Unknown'' level, and count features are filled with zero, which correctly represents an absence of events.

The number of records in the observation table is in millions, so the ingestion is done in parts. This step produces normalized data that includes patients, encounters, conditions, medications, procedures, observations, care plans and providers.

\subsection{Stage 2: Feature Engineering and Clustering}
For clustering, we need a table in which every row represents a patient and every column represents a feature. Most effort here goes into creating such a table, since the clinically relevant information is scattered across the database tables and not directly available for modeling. For each feature, a separate SQL aggregation query is run.

We start off with a large collection of potential features and end up selecting a few non-overlapping ones later. These features are grouped according to type in Table~\ref{tab:features}. The numeric features are relatively simple (counts, lengths and ratios). On the other hand, there is more domain knowledge embedded within the categorical features because we do not rely on raw codes, which are derived into clinically meaningful tiers through rule-based SQL. Disease features such as diabetes, hypertension, heart disease, cancer and kidney disease are defined based on keyword matching against condition descriptions, specifically excluding ``prediabetes'', ``suspected'' or ``family history'' when the feature would imply otherwise. Care behavior features include a tier of thirty-day readmissions generated using window functions on consecutive visits, a flag for those who make visits to the ED in out-of-office hours, their drug use pattern which integrates their diagnostic information with the results of the AUDIT-C and DAST-10 tests, and a tier for different pain management approaches.

\begin{table}[t]
\caption{Engineered patient features, grouped by kind. The full catalogue is computed offline. Clustering uses a non-redundant subset.}
\label{tab:features}
\centering
\footnotesize
\begin{tabular}{@{}p{0.26\linewidth}p{0.66\linewidth}@{}}
\toprule
Group & Features \\
\midrule
Numerical & age, active conditions, total medications, total procedures, distinct providers seen, lifetime medication days, encounter duration hours, preventive-to-reactive ratio \\[2pt]
Demographic & gender, marital status, race, ethnicity \\[2pt]
Disease flags & diabetes, hypertension, heart disease, cancer, kidney disease, mental health, respiratory profile \\[2pt]
Binned vitals & BMI class, blood-pressure risk, glucose risk (each with an explicit ``Unknown'' level) \\[2pt]
Care-pattern tiers & dominant care plan (severity-ranked), condition burden, medication intensity, polypharmacy, 30-day readmission tier, out-of-hours flyer \\[2pt]
Risk profiles & substance-use profile, pain-management tier, cholesterol risk, pregnancy history, smoking status \\
\bottomrule
\end{tabular}
\end{table}

The selected engineering methods include the following. Some features are heavily skewed, so we apply a natural log transform and then min-max scaling. For the physiological vital sign features such as BP, Glucose, we assume that they may not be available, hence categorizing them along with unknowns instead of replacing the missing values with mean values as part of retaining the knowledge that the particular feature was never recorded. The most critical care plan is not the most common care plan but the care plan that is most severe. We arrange the care plans per patients according to specific severity rankings (oncology and palliative, maternal, trauma, respiratory and cardio-metabolic), then choose the most severe care plan. So a cancer patient is assigned the cancer care plan even though their most common care plan can be a checkup. We do not transform categorical features through one hot encoding because K-Prototypes can work with them directly. There is also the clinical severity score (first principal component of log-transformed counts of procedures/medications/conditions).

K-Prototypes algorithm \cite{huang1998} is an extension of K-Means algorithm for mixed data. Distance is defined by the sum of squares of the difference of numerical features and the number of mismatches in categorical features. The weighted sum using weight $\gamma$ helps to control the relative significance of the categorical features. We use the Huang initialisation and select the value of $\gamma$ automatically using the library.

Prior to the model fitting, we look for any statistical outliers based on a z-score criterion for the numerical variables after log-transforming the heavily skewed counts, such that the tail distribution itself does not lead to the detection of an outlier. If a z-score greater than six is detected for any of the variables, the corresponding patient is considered an outlier and there are 289 such patients that constitute around 1\% of the dataset. None of the outliers are removed from the analysis. Instead, they receive a unique label of $-1$ and are included in the analysis with the other patients without being used in the K-Prototypes cluster fitting, i.e., none of the clusters are generated for them. This is done because their values would skew the cluster generation, but they are clinically important since they represent the most severe cases like cancer, dialysis, and heart disease.

\subsection{Stage 3: Clinical Sequence Mining}

In the third phase, we study the order of clinical events. One ordering sequence is defined for each patient, i.e., the ordering sequence of all the events occurring in a certain timeline, arranged according to the dates. Each event is prefixed with a type character such that the source table for the particular event is known (Condition is labeled C, Medication M, Procedure P, and Care plan CP). If more than one event occurs on a particular day, then the ordering is based on a clinical order (C-M-P-CP).

Cleaning comes before mining, where we remove all events except for those occurring during the last two years per patient, exclude some purely administrative events and psychosocial ``status'' events added to almost any encounter by Synthea, merge immediate duplicates of events, and restrict the length of the event timeline, allowing us to perform mining very efficiently. Mining of sequences is done using the PrefixSpan algorithm \cite{pei2001}, a classical approach that builds on finding closed frequent sequences by projecting the data onto all the already discovered prefixes. We ask PrefixSpan to generate closed sequences of length ranging from three to six events having a support of at least one percent of patients, and we mine more than needed for selection purposes.

\begin{figure}[t]
\centering
\includegraphics[width=0.8\linewidth]{figs/sequence_mining_pipeline.png}
\caption{The sequence-mining pipeline}
\label{fig:pipeline}
\end{figure}

The key step happens around PrefixSpan in Figure~\ref{fig:pipeline}. When ranking the patterns by their supports only, the best patterns are going to be very frequent but not useful sequences, such as regular dentistry cleaning actions which everyone undergoes. To solve this, we rank the patterns using Geometric Mean Lift (GML). Given a pattern $p$ containing $k$ elements with its support being $s(p)$ among $n$ patients and item supports being $s(i)$, we compute
\begin{equation}
\mathrm{GML}(p) \;=\; \frac{s(p)/n}{\left(\prod_{i \in p} s(i)/n \right)^{1/k}} .
\end{equation}
It is just another step in applying the notion of lift from two items to $k$ items based on the geometric mean calculation. It can be calculated in log scale to prevent overflow. The higher GML value is, the more probable it is that the pattern would appear in practice more often than its component frequencies' product.

Two other concepts have been proposed. Firstly, an additional bonus is assigned to the overall score as multiplication by a multiplier of $(1 + 0.3\, (t-1))$ where $t$ stands for the total number of different types of events in the pattern. In this way, we give preference to the actual trajectory of a patient which includes diseases, drugs, and surgeries to a technical pattern of three drug applications. Secondly, certain diversity criteria are used. We filter out patterns including triplets of the same type of events, patterns with very similar items according to the Jaccard distance, and limit the number of patterns including the same beginning or ending events or dominated by dentistry events. Also, some slots are reserved for fully representative patterns that include all four types of events.

Finally, all these processes are run for each scope once. Mining is performed both on the global scope and on the individual clusters, including the outlier group denoted $-1$. Patterns discovered in the global scope are ignored while mining the individual clusters so as to keep only those patterns that define the clusters. Further, each of the patterns retained has their patient ID attached to them using the intersecting posting list of the least frequent items.

\subsection{AI Interpretation Layer}

Finally, this layer makes the output understandable. We use just one LLM using the API of OpenAI. However, we do this wisely: the statistics are never computed by the machine alone. The programming code computes everything before providing these static values to the machine in its prompt, asking it to verbalize about no other data but the provided numbers. Such an approach is known as ``verbalizer'' since it guarantees that facts are described in the text.

The key point is that clinical features are given as auxiliary data within the prompt. Clusters' prompts are constructed out of actual figures: population averages, cluster-specific values with an annotation ``compared to population average'', prevalence rates of various conditions, averages for vital indicators, dominant care plan category, and prevalent diseases, therapies, and interventions. In turn, the model is supposed to generate strict JSON including the risk level, clinical reasoning for doctors, statement for patients, key characteristics of a certain cohort or sequence, and suggestions as to what actions should be taken. Sequences interpretation is done analogically by providing mined patterns accompanied by aggregated patient profile (age, gender, BMI, disease prevalence, dominant care plan), and further decoding the prefixes and temporal arrows into reasonable care pathway.

All AI outputs are cached as a JSON value together with their corresponding rows, so the model is run only once per row.

\subsection{Natural-Language-to-SQL Agent}

In addition to strict interpretations, HospitAI allows users to ask questions on top of the available data. A question-answer agent was developed, which is a state machine based on LangGraph, shown in Figure~\ref{fig:agent}. Each state accomplishes one task alone.

\begin{figure}[t]
\centering
\includegraphics[width=1\linewidth]{figs/agent_diagram.png}
\caption{The natural-language-to-SQL agent}
\label{fig:agent}
\end{figure}

In the initial phase the question is classified and rephrased. The agent determines whether the question requires the database, whether it is a generic medical question that can be answered directly, or whether it is irrelevant and must be rejected.

In case the database needs to be accessed, the query generation procedure is performed. Generation is based on the embedding of the database schema and a set of rules. Generated queries are validated in two layers. The first, deterministic layer ensures the query is a single read-only statement and blocks potentially harmful keywords via regex validation. The second layer validates the query semantics, i.e. its ability to answer the question with the appropriate tables and filters. If the query fails either layer, the error message is passed back to the agent to regenerate the query until an attempt limit is reached. Database execution errors are treated the same way.

The agent is also role-aware. If the request comes from a doctor, an additional scope condition restricts every search to the records of that doctor's own patients via their provider id, and any request without this condition is rejected, preventing a doctor from seeing other doctors' patients.
\section{Results}

\label{sec:results}

The effectiveness of the proposed approach is illustrated using the Synthea dataset, which is introduced in Section~\ref{sec:system}. Out of 22{,}284 patients, only 21{,}995 patients have been clustered, and all outliers have been retained in a different class labeled as $-1$.

\subsection{Clustering}

To discover the optimal configuration, we performed a large benchmarking procedure on 628 configurations. Configuration means feature subset selection (we tried more than sixty feature subsets), cluster number $k$, scaler used, categorical weight $\gamma$, random state and outlier threshold. For benchmarking, we calculated a score based on three metrics. First of all, the metrics are min-max normalized individually within the scope of all 628 configurations, and the score is then their linear combination with weights that reflect our confidence in them ($0.5$ for Gower silhouette, $0.3$ for Calinski-Harabasz index and $0.2$ for Davies-Bouldin index, inverted because the higher the score the worse the partition is). We chose the best clinically interesting configuration per $k$ value (we tried four, five and six), and from these three configurations we tried once again to optimize the hyperparameters and feature subset used to choose the best local configuration.

\begin{table*}[t]

\caption{Representative clustering configuration for each cluster number (feature vector based on the corresponding version id). CH index and DB index are computed from the FAMD embedding in the mixed score framework. The highest value obtained for each criterion is marked in bold font. The selected cluster number is $k=4$ (v55).}

\label{tab:clusters}

\centering

\footnotesize

\begin{tabular}{@{}l c c l c r c@{}}

\toprule

Config & $k$ & Gower & Per-cluster Gower & DB & CH & Score \\

\midrule

v55 (final) & 4 & \textbf{0.350} & 0.42, 0.20, 0.32, 0.40 & \textbf{0.72} & \textbf{26813} & \textbf{0.82} \\

v40 & 5 & 0.324 & 0.28, 0.22, 0.31, 0.19, 0.48 & 0.87 & 22547 & 0.74 \\

v60 & 6 & 0.287 & 0.17, 0.13, 0.47, 0.21, 0.38, 0.27 & 0.95 & 22273 & 0.70 \\

\bottomrule

\end{tabular}

\end{table*}

Table~\ref{tab:clusters} reports a representative configuration for each $k$. The choice of $k$ is made to be 4, with feature combination v55. It scores best overall (0.82) among the clinically meaningful runs, and does so across all individual metrics: the largest Gower silhouette, the largest Calinski-Harabasz score, and the smallest Davies-Bouldin score. We still eliminate one family of runs, which scores even better on paper. It achieves a score of 0.87 by segregating patients according to gender, but since the separation is meaningless clinically, we exclude them. Among the four clusters in the optimal run, three get a Gower silhouette per cluster greater than 0.30, and only one gets close to 0.20, compared with the results of five and six clusters where one and two clusters respectively have values less than 0.20. Since clusters will be used for the scope of the sequence mining analysis, we prefer to have the smallest $k$ possible while still having all the clusters be cohesive. The important features are the number of active conditions and patient age numerically, and condition burden, gender, diabetes, and heart disease categorically.

Figure~\ref{fig:heatmap} shows the per-cluster means of the numeric variables represented as a heatmap. There is a clear morbidity gradient. That is, cluster 0 represents young people with low morbidity, while the rest of the clusters contain a progressively greater number of conditions. By interpreting all four clusters in conjunction with the categorical features, we see four distinct patient profiles, which the AI platform automatically labels as: Younger Injury Care (young, low morbidity), Pregnancy Focused Adults (reproductive women), Older High Risk (hypertensive polypharmacy adults), and Kids and Well Visits (pediatric/acute patients). These titles, and the description, insight, and recommendation attached to each cluster, are produced by the LLM at each run and are therefore not deterministic.

\begin{figure}[t]

\centering

\includegraphics[width=0.95\linewidth]{figs/cluster_heatmap.png}

\caption{Z-scored numerical feature means per cluster ($k=4$). Warm colours are above the population average, cool colours below it.}

\label{fig:heatmap}

\end{figure}

\subsection{Sequence Mining}

The clearest sign of our pipeline's impact is the change in the mined patterns' properties as each new version is introduced. In Table~\ref{tab:ablation}, we present two sample patterns of at least four events from every one of these versions. The initial versions (v0-v2) rank patterns solely based on support, and therefore, the beginning of the ranking consists of common yet meaningless sequences, which consist mainly of similar dental cleanings that about 93\% of the patients undergo. With v3 and onwards, geometric mean lift-based ranking along with diversity filters were applied, and at the top of the ranking, clinically consistent care pathways can be seen. Version v4 also filters out consecutive occurrence of similar events, which is why the long sequence of procedure events related to a diabetic eye exam no longer appears from v4 onwards. At v6, when social status events are removed and reserved space for full spectrum patterns was set aside, diverse combinations of conditions, care plans, medicines and procedures appear despite their low support scores.

\begin{table}[t]

\caption{How the mined patterns evolve across miner versions}

\label{tab:ablation}

\centering

\scriptsize

\begin{tabular}{@{}l p{0.72\linewidth} r@{}}

\toprule

Ver. & Example pattern (length $\ge 4$) & Sup. \\

\midrule

v0 & dental consultation $\rightarrow$ dental care $\rightarrow$ supragingival scaling $\rightarrow$ subgingival scaling & 93\% \\

v0 & supragingival scaling $\rightarrow$ subgingival scaling $\rightarrow$ gingivae exam $\rightarrow$ oral health education & 93\% \\

\addlinespace

v1 & supragingival scaling $\rightarrow$ subgingival scaling $\rightarrow$ gingivae exam $\rightarrow$ oral health education & 63\% \\

v1 & dental care $\rightarrow$ supragingival scaling $\rightarrow$ subgingival scaling $\rightarrow$ oral health education & 63\% \\

\addlinespace

v2 & dental consultation $\rightarrow$ dental care $\rightarrow$ subgingival scaling $\rightarrow$ gingivae exam $\rightarrow$ oral health education & 62\% \\

v2 & dental referral $\rightarrow$ dental consultation $\rightarrow$ dental care $\rightarrow$ subgingival scaling $\rightarrow$ oral health education & 57\% \\

\midrule

v3 & tropicamide eye drops $\rightarrow$ gonioscopy $\rightarrow$ fundoscopy $\rightarrow$ diabetic eye exam & 3\% \\

v3 & furosemide $\rightarrow$ oxygen by mask $\rightarrow$ prone position $\rightarrow$ furosemide & 2\% \\

\addlinespace

v4 & acute bronchitis $\rightarrow$ respiratory therapy $\rightarrow$ respiratory function test $\rightarrow$ acetaminophen & 3\% \\

v4 & salmeterol-fluticasone inhaler $\rightarrow$ albuterol $\rightarrow$ spirometry $\rightarrow$ depression screening & 2\% \\

\addlinespace

v5 & subgingival scaling $\rightarrow$ gingivae exam $\rightarrow$ sodium fluoride gel $\rightarrow$ fluoride treatment & 35\% \\

v5 & dental bitewing X-ray $\rightarrow$ gingivae exam $\rightarrow$ sodium fluoride gel $\rightarrow$ fluoride treatment & 19\% \\

\addlinespace

v6 & acute bronchitis $\rightarrow$ respiratory therapy $\rightarrow$ respiratory function test $\rightarrow$ acetaminophen & 3\% \\

v6 & chest X-ray $\rightarrow$ furosemide $\rightarrow$ oxygen by mask $\rightarrow$ prone position $\rightarrow$ furosemide & 1\% \\

\bottomrule

\end{tabular}

\end{table}

Frequent pattern mining from each scope individually yields 175-252 frequent patterns per cluster. Examples of representative frequent patterns are listed in Table~\ref{tab:patterns}. Each cluster has a unique signal dominating its set of frequent patterns. For example, in Cluster~0, the dominant signal relates to acute care episodes. In Cluster~2, the chronic disease management includes antihypertensive drugs combined with screenings, or furosemide along with oxygen therapy to manage fluid overload. The most valuable benefit is obtained from the outlier scope. Since we mine each individual scope independently, high-cost combinations like chemotherapy and cisplatin leading to hospice care, or dialysis and epoetin, will dominate the results. Otherwise, these pathways fall under the global support threshold and would be missed when mining all cases at once.

\begin{table}[t]

\caption{Representative mined pathways by scope (ranked by GML).}

\label{tab:patterns}

\centering

\footnotesize

\begin{tabular}{@{}p{0.16\linewidth}p{0.74\linewidth}@{}}

\toprule

Scope & Pattern \\

\midrule

Cluster 0 & Fracture of bone $\rightarrow$ Fracture care $\rightarrow$ Bone immobilization \\

Cluster 1 & Normal pregnancy $\rightarrow$ Pregnancy test $\rightarrow$ Ultrasound for fetal viability \\

Cluster 2 & Furosemide $\rightarrow$ Oxygen by mask $\rightarrow$ Prone position $\rightarrow$ Furosemide \\

Cluster 3 & Acute allergic reaction $\rightarrow$ Epinephrine injection $\rightarrow$ Prednisone \\

Complex High Need ($-1$) & Chemotherapy + radiation $\rightarrow$ Cisplatin $\rightarrow$ Hospice care \\

\bottomrule

\end{tabular}

\end{table}

\subsection{AI Layer}

Upon running the K-Prototypes script, the AI creates a short card per cluster based on the statistics generated beforehand and determines the cluster title, its colour theme that will be used in all mentions of the cluster on the platform, its description, an insight and a recommendation for each. Figure~\ref{fig:platform_clusters} shows what these cards look like in the dashboard. Since the model uses only the precomputed values, the card matches the heatmap and the prevalence rates by assigning high risk and a cardiology-related recommendation to the high-burden cluster, and a healthy, preventive profile otherwise.

\begin{figure}[t]

\centering

\includegraphics[width=\linewidth]{figs/platform_clusters.png}

\caption{The clusters as shown in the platform. The title, colour, description, insight and recommendation of each card are produced automatically by the AI when the clustering script runs.}

\label{fig:platform_clusters}

\end{figure}

Beyond the card information, the platform generates a more detailed report about every identified cluster or mined pattern. It serves two purposes: a clinical explanation for clinicians, and a simpler, layman's summary for someone without a medical background. This kind of explanation is depicted in Figure~\ref{fig:platform_interpretation}.

\begin{figure}[t]

\centering

\includegraphics[width=\linewidth]{figs/platform_interpretation.png}

\caption{A cluster AI interpretation in the front end, written for both a clinician and a reader without medical training.}

\label{fig:platform_interpretation}

\end{figure}

Finally, the natural-language-to-SQL agent lets a user ask free questions. We tested it on questions of different difficulty, including non-trivial ones that need several joins. Table~\ref{tab:agent_example} lists two such questions: the first joins several clinical tables, the second joins the cluster and the sequence-mining tables. The answers are correct almost every time, and the guard layer blocked every attempt to run a non-\texttt{SELECT} statement. When the agent cannot build a valid query or finds no matching data, it plainly says it could not complete the task instead of inventing an answer.

\begin{table}[t]

\caption{Two example exchanges with the natural-language-to-SQL agent. The first needs several joins across the clinical tables, while the second joins the cluster and the sequence-mining tables.}

\label{tab:agent_example}

\centering

\footnotesize

\begin{tabular}{@{}p{0.48\linewidth}p{0.44\linewidth}@{}}

\toprule

Question & Agent answer \\

\midrule

How many distinct female patients over 60 with a diabetes diagnosis were also prescribed metformin and had at least one emergency encounter? & There were 144 distinct female patients over 60 with a diabetes diagnosis who were prescribed metformin and had at least one emergency encounter. \\[4pt]

Which mined care pathway is followed by the most patients in cluster~2, and how many of them follow it? & In cluster~2, the most common mined pathway is ``M:sodium fluoride 0.0272 MG/MG Oral Gel $\rightarrow$ P:Dental fluoride treatment $\rightarrow$ P:Oral health education'', followed by 1{,}265 patients. \\

\bottomrule

\end{tabular}

\end{table}

A full quantitative evaluation of the agent and of the generated summaries is left for future work, as discussed next.


\section{Discussion and Conclusion}
\label{sec:conclusion}

We have introduced HospitAI, which brings together three analytical stages in hospitals based on data analysis: clustering patients using mixed-type clustering, sequence mining, and AI interpretation. The clustering produces a small set of understandable clusters of patients. Sequence mining, enhanced by geometric mean lift scoring and diverse filtering, provides clinically relevant paths rather than statistically frequent ones. Language generation converts numerical representations into short summaries and allows natural language querying.

There are also several limitations. The data is artificial, meaning the trends represent modeling in Synthea and do not depict an actual hospital's processes. We also have not conducted any form of clinical validation with medical practitioners, and the mined trends were evaluated based on qualitative evaluation rather than quantitatively compared to other miners with a benchmark for the same data set. The system depends on one external language model, which introduces challenges of expense and repeatability and poses a problem of governance that needs to be resolved before deployment at a hospital. The filter we used is tuned for the data, implying that some of the cleaning process is unique to this dataset.

Several next steps follow. First, we need to validate our approach on a real de-identified dataset and have clinicians evaluate our mined pathways. Second, a better evaluation of the sequence patterns would be beneficial, for instance by comparison with existing clinical guidelines or against other mining algorithms. Finally, a detailed analysis of the language layer is needed, particularly the quality of the SQL agent and of the produced summaries. Replacing the dataset-specific filtering with a more generic one would also make the pipeline more transferable. With appropriate preprocessing and postprocessing, off-the-shelf algorithms already give a hospital an affordable platform built on data it should already have.


\bibliographystyle{IEEEtran}
\bibliography{references}

\end{document}
